\chapter{多次调和函数}
\section{上半连续函数}

\begin{definition}
    $\Omega \subset \rr^n$ 区域. 称函数 $u:\Omega \to [-\infty ,\infty )$ 为上半连续的, 若对 $\forall c \in \rr$, 有 $\{u<c\} \subset \Omega $ 为开集($\iff \limsup_{z \to a}u(z)=u(a), \forall a \in \Omega $ ).  
    记 $\usc(\Omega )$ 为 $\Omega $ 上的上半连续函数的全体(upper semi-continuous).     
\end{definition}
\begin{remark}
    \textcolor{purple}{若按微积分中的上极限定义, 趋于 $a$ 的点列不含 $a$ 本身, 则上述定义中的 $\limsup_{z \to a}u(z) = u(a)$ 应改为 $\limsup_{z \to a}u(z)\leq u(a)$; 若趋于$a$的点列可以包含 $a$, 则自然有 $\limsup_{z \to a}u(z)\geq u(a)$; 上半连续也等价于: 对任意的 $x \in \Omega, f(x)\neq -\infty  $以及 任意的 $\epsilon >0$, 存在邻域 $U \ni x$ s.t. $f(y)<f(x)+\epsilon, \forall y \in U$.    }
\end{remark}
\textbf{基本性质}:
\begin{enumerate}[(1)]
    \item $u \in \usc(\Omega )$ $\Rightarrow $ 对任意的紧集 $K \subset \Omega $, $u$ 在 $K$ 上可取到最大值;
    \item (\textbf{Dini}): $\{u_{j}\} \subset \usc(\Omega ), u_{j} \searrow 0$ $\Rightarrow $ 对任意的紧集 $K \subset \Omega $ 有 $u_{j} \rightrightarrows 0$ 于 $K$;
    \item $u \in \usc(\Omega )$ 且有上界 $\Rightarrow \exists u_{j} \in C(\Omega )$ s.t. $u_{j} \searrow u$.             
\end{enumerate}
\begin{proof}
    (1): 令 $U_{j}:=\{u <j\}$, 则 $U_j \nearrow$,  且为开集, 且 $\cup _{j}U_j \supset K$,
    于是 $\exists j_0$ s.t. $K \subset U_{j_0}$. $\Rightarrow S:=\sup_K u \leq j_0<\infty $, 取 $\{x_j\} \subset K, u(x_j) \to S$ $\Rightarrow \exists $ 子列 $x_{j_k} \to x_0 \in K$.
    则 $S=\lim u(x_{j_k}) \leq u(x_0)$ (上半连续性) $\leq S$, 所以 $u(x_0)=S$.   
    
    (2): $\forall \epsilon >0$, 取 $U_j:=\{u_j <\epsilon \}$, 则 $U_j$ 为开集 且 $U_j \nearrow$. 因为 $u_j \searrow 0,$ 所以 $K \subset \cup _j U_j$. 于是存在 $j_0$ s.t. $K \subset U_{j} \forall j \geq j_0$, 即$0 \leq u_j(x) <\epsilon, \forall x \in K, \forall j \geq j_0 $. 
    
    (3): 令 $u_j:=\sup_{y \in \Omega } \{u(y)-j |x-y|\}$. \\
    $u_j \in C(\Omega )$:  设 $x,y \in \Omega .$ 则$ \forall z \in \Omega $ 有 
    $$ \begin{aligned}
       & u(z)-j|x-z| \leq u(z)-j |y-z|+j |x-y| \\
       \Rightarrow \quad & u_{j}(x) \leq u_j (y)+j |x-y| \\
       \Rightarrow \quad & u_{j}(y) \leq u_j (x)+j |x-y| \\
       \Rightarrow \quad& | u_j(x)-u_j(y) | \leq j|x-y|.
    \end{aligned}  $$
  另外, $u_j \searrow $ 且 $u_j(x) \geq u(x)$. 设 $u \leq M$, $u(x) \neq -\infty $. 则 $\forall \epsilon >0, \exists \delta >0$, s.t. 
  $$ u(y)<u(x)+\epsilon , \quad \forall y \in B_{\delta }(x). $$     
  当 $|x-y| \geq \delta $时, 
  $$ u(y)-j|x-y| \leq M-j \delta  < u(x), \quad (j \geq j_0(x,\epsilon )>>1) $$
  $$ \Rightarrow \quad u(x) \leq u_j(x) \leq u(x)+\epsilon  $$
   $$ \Rightarrow \quad u_j (x) \to u(x).$$
若 $u(x)=-\infty $. 则 $\forall N>0, \exists \delta >0$, s.t. 
$$ u(y)<-N, \quad \forall y \in B_{\delta }(x). $$  
当 $|x-y| \geq \delta $ 时, 
$$ u(y)-j|x-y| \leq M -j \delta  \leq -N, \quad (j \geq j_0>>1) $$
$$\Rightarrow \quad u_j(x ) \leq -N  \quad \Rightarrow \quad u_j(x) \to -\infty =u(x).$$
~~
\end{proof}

\begin{definition}
    设函数 $u:\Omega \to [-\infty ,+\infty )$, 局部有\textcolor{purple}{(上)}界. 称 $u^{*}(x):=\limsup _{y \to x}u(y)$ 为 $u$ 的上半连续化.   
\end{definition}
(4)(\textbf{续}) $u \in \usc(\Omega ) \iff u=u^{*}$. 
\begin{proof}
    $\Rightarrow $: 设 $u(x) \neq -\infty $, $\forall \epsilon >0, \exists \delta >0$, s.t. $u(y) <u(x)+\epsilon , \forall y \in B_{\delta }(x)$. $\Rightarrow u^*(x) \leq u(x)+\epsilon $.  由 $\epsilon $ 的任意性知, $u^*(x) \leq u(x)$. 又 $u^*(x) \geq u(x) \Rightarrow u^*(x)=u(x)$.
    
    $\Leftarrow $: 只需证 $u^*$ 为上半连续. 
    设 $u^*(x) \neq -\infty $ (\textcolor{red}{$-\infty $ 的情况作为练习}).   
    $\forall \epsilon >0, \exists \delta >0,$ s.t. 
    $$ u(y)<u^*(x)+\frac{\epsilon }{2}, \quad \forall y \in B_{\delta  }(x) .$$ 
    设 $z:|z-x|< \frac{\delta }{2}$, 取 $y \in B_{\delta }(x)$ s.t. 
    $$ u^*(z) < u(y)+\frac{\epsilon }{2} < u^*(x)+\epsilon  $$
    $\Rightarrow $ $u^*$ 在 $x$ 处USC.     
\end{proof}




\section{次调和函数}

\paragraph{A. 调和函数}
\begin{definition}
    设 $\Omega \subset \cc$ 为区域, $u \in C(\Omega , \rr)$, 若 $a \in \Omega $ 以及 $\Delta (a,r) \subset \subset \Omega $: 
    $$ u(a)=\frac{1}{2 \pi } \int_{0}^{2 \pi } u(a+r e^{i \theta }) d \theta \quad (\text{均值性质}) ,$$    
    则 $u$ 为 $\Omega $ 上的调和函数, 其全体记为 $H(\Omega )$.   
\end{definition}

\noindent\textbf{基本性质}:
\begin{enumerate}
  \item 最大值原理：
  \(
    u \in H(\Omega) \text{ 不能在内部取最大值}.
  \);
  \item
  \(
    u \in H(\Omega)
    \Longleftrightarrow
    u \in C^{\infty }(\Omega) \text{ 且 } \,\Delta u = 0;
  \)
  \item Dirichlet 问题：
  设
  \(
    f \in C(\partial \Delta(a,r), \mathbb{R}),
  \)
  则
  \[
    u(z)
    :=
    \frac{1}{2\pi}
    \int_0^{2\pi}
    f(a+re^{i\theta})
    \frac{r^2 - |z-a|^2}{|re^{i\theta}-(z-a)|^2}
    \, d\theta , \quad (\text{Poisson 积分})
  \]
满足
  \(
    u \in H(\Delta(a,r)) \cap C(\overline{\Delta(a,r)}),
    \,
    u|_{\partial \Delta(a,r)} = f 
  \);
  \item
  \(
    u \in H(\Omega)
    \Longleftrightarrow
    u \text{ 局部表示为一个全纯函数的实部}
  \)
   ($\Rightarrow  u$ 为实解析); 
  \item Harnack 定理：
  若
  \(
    u_j \in H(\Omega)
  \)
  且
  \(
    u_j \to  u (\text{内闭匀敛}),
  \)
  则
  \(
    u \in H(\Omega).
  \)
\end{enumerate}

\paragraph{B. 次调和函数}

\begin{definition}
设 \(\Omega \subset \mathbb{C}\) 为区域， \(u \in \operatorname{USC}(\Omega)\)。
若对任意 \(a \in \Omega\) 以及任意
\[
  \Delta(a,r) \subset \subset  \Omega,
\]
有
\[
  u(a)
  \leq
  \frac{1}{2\pi}
  \int_0^{2\pi}
  u(a+re^{i\theta})\, d\theta \quad (\text{次均值性}),
\]
则称 \(u\) 在 \(\Omega\) 上次调和，记为
\(
  u \in \sh(\Omega).
\)
\end{definition}

\begin{example}
若 \(u \in H(\Omega)\)，则
\(
  |u|^\alpha \in \sh(\Omega),
  \forall \alpha \geq 1.
\)
\end{example}
\begin{proof}
    $$ \begin{aligned}
        \left|u(a)\right|
\leq &
\frac{1}{2\pi}
\left|
\int_0^{2\pi} u(a+re^{i\theta})\, d\theta
\right|\\
  \leq &
\frac{1}{2\pi}
\left(
\int_0^{2\pi}
\left|u(a+re^{i\theta})\right|^\alpha
\, d\theta
\right)^{1/\alpha}
\cdot
(2\pi)^{1-1/\alpha}.     
    \end{aligned}  $$
    ~~
\end{proof}

\subsubsection{基本性质}
\begin{enumerate}[(1)]
  \item $u_1,u_2 \in \sh(\Omega ), \alpha _1,\alpha _2 \geq 0$ $\Rightarrow \alpha _1 u_1+\alpha _2 u_2 \in \sh(\Omega )$ (\textcolor{red}{锥});
  \item   $u_1,u_2 \in \sh(\Omega )$ $\Rightarrow \max\{u_1,u_2\} \in \sh(\Omega )$;
  \item $\mathcal{F} \subset \sh(\Omega )$, 若 $v:=\sup_{u \in \mathcal{F}}u$ 局部有上界 $\Rightarrow v^{*} \in \sh(\Omega )$. 
\end{enumerate}
\begin{proof}
  (3): 设 $\Delta (a,r) \subset \subset \Omega $: 
  $$ u(a) \leq  \frac{1}{2 \pi } \int_{0}^{2 \pi } u(a+r e^{i \theta }) d \theta \leq \frac{1}{2 \pi } \int_{0}^{2 \pi } v^{*}(a+r e^{i \theta }) d \theta ,$$ 
  于是
  $$  \quad v(a )\leq \frac{1}{2 \pi } \int_{0}^{2 \pi } v^{*}(a+r e^{i \theta }) d \theta .$$
  取 $a_j \to a$, s.t. $v(a_j) \to v^*(a)$, 则 
  $$ \begin{aligned}
    v^*(a)=\lim v(a_j) \leq & \limsup _{j \to \infty } \frac{1}{2 \pi } \int_{0}^{2 \pi } v^{*}(a_j+r e^{i \theta }) d \theta \\
    \textcolor{red}{(\text{Fatou:利用} c-v^* \geq 0)} \leq &  \frac{1}{2 \pi } \int_{0}^{2 \pi } \limsup _{j \to \infty } v^{*}(a_j+r e^{i \theta }) d \theta \\
    =& \frac{1}{2 \pi } \int_{0}^{2 \pi } v^{*}(a+r e^{i \theta }) d \theta.
  \end{aligned}  $$  
  ~~
\end{proof}

\begin{example}\label{ex2.2}
  设 $\Omega \subsetneq \cc$ 为一区域, 令 $\delta (z):=d(z, \partial \Omega )$. 则 $-\log \delta  \in \sh (\Omega )$.
  
  \textbf{解}: $$\delta (z)=\inf_{w \in \partial \Omega }|z-w| \Rightarrow -\log \delta (z)=\sup_{w \in \partial \Omega }\log \frac{1}{|z-w|}.$$
  因为对每个 $w \in \partial \Omega , \log \frac{1}{|z-w|} \in H(\Omega )$, 且 $-\log \delta \in C(\Omega )$, 所以由基本性质(3) 得   $-\log \delta  \in \sh (\Omega )$.
\end{example}
\begin{remark}
  $\Omega $ 凸 $\Rightarrow -\delta $ 凸函数(\textcolor{red}{反之是否成立?}).    
\end{remark}
\begin{remark}
  $\log |f| \in H(\Omega )$: 
  $$ \log|f|=\frac{1}{2}\underbrace{\log f}_{\text{全纯}} + \frac{1}{2}\underbrace{\overline{\log f}}_{\text{反全纯}} ,$$
  而 $\Delta =4 \partial \bar \partial $. 
\end{remark}


\noindent (4)
$\{u_j\} \subset \sh (\Omega )$ 且 $u_j \searrow u, \Rightarrow u \in \sh(\Omega )$.
  \begin{proof}
    $u \in \usc:$
    $$ \begin{aligned}
     & \limsup _{z \to a} u(z) \leq \limsup _{z \to a} u_j(z )=u_j(a) \to u(a)\\
     & \Rightarrow \quad \limsup _{z \to a} u(z)=u(a) .
    \end{aligned}  $$ 
    设 $\Delta (a,r) \subset \subset \Omega $, 则 
    $$ u_j(a) \leq \frac{1}{2 \pi }\int_{0}^{2 \pi } u_j(a+r e^{i \theta }) d \theta  .$$ 
    令 $j \to \infty $, 由 Levi定理以及 $u_j$ \textcolor{purple}{(上)}有界得 
    $$ u(a) \leq \frac{1}{2 \pi }\int_{0}^{2 \pi } u(a+r e^{i \theta }) d \theta  .$$ 
    ~~
  \end{proof}

\noindent (5) $\lambda$ 为 $\rr$ 上的凸增函数, $u \in \sh(\Omega ) \Rightarrow \lambda \circ u \in \sh(\Omega )$.  
\begin{proof}
  因为 $\lambda $ 连续, 所以 $\lambda \circ u \in \usc(\Omega )$.   因为 $\lambda $ 凸, 所以 $\forall (x_0,\lambda (x_0))$, 存在直线 $y=\lambda (x_0)+k(x-x_0)$, s.t. 曲线 $y=\lambda (x)$位于直线上方. 
  $$ \Rightarrow \lambda (x) \geq \lambda (x_0)+k(x-x_0) .$$
  设 $\Delta (a,r) \subset \subset \Omega $. 取 $x=u(a+r e^{i \theta })$ 代入上式, 再积分: 
  $$ \frac{1}{2 \pi }\int_{0}^{2 \pi } \lambda \circ u(a+r e^{i \theta }) d \theta  \geq \lambda (x_0)+ k \big\{\frac{1}{2 \pi } \int_{0}^{2 \pi } u(a+r e^{i \theta }) d \theta -x_0 \big\} .$$      
  取 $x_0= \frac{1}{2 \pi } \int_{0}^{2 \pi } u(a+r e^{i \theta }) d \theta $ $(\geq u(a))$, 
  $$ \Rightarrow \quad  \frac{1}{2 \pi }\int_{0}^{2 \pi } \lambda \circ u(a+r e^{i \theta }) d \theta \geq \lambda \circ u(a).$$ 
  ~~
\end{proof}


\begin{definition}
  设 $u \in \usc(\Omega )$, 若 $\forall a \in \Omega $, 存在邻域 $U_a \ni a$, s.t. $u \in \sh(U_a)$, 则称 $u$ 为局部次调和的, 记全体为 $\sh_{loc}(\Omega )$.      
\end{definition}

\noindent (6)
若 $u \in \sh_{loc}(\Omega )$, 在 $\Omega $ 内部达到最大值, 则 $u$ 为常数. ($\Rightarrow $ 最大模原理: $|f|^2=(\Re f)^2+(\Im f)^2 \in \sh, (f \in \mathcal{O})$ .)   
\begin{proof}
  设 $a \in \Omega $, s.t. $u(a)=\max_{\Omega }u$.
  取 $r_a >0$, s.t. $u\in \sh(\Delta (a,r_a))$.
  下证 $u|_{\partial \Delta (a,r)}=u(a), \forall a<r<r_a$:   
  若 存在 $\theta _0$ s.t. 
  $$ u(a+re^{i \theta _0})<u(a) ,$$ 
  由 $\usc$ 性, $\exists \epsilon _0 <<1$ s.t. 
  $$ \frac{1}{2 \pi } \int_{\theta _0 -\epsilon }^{\theta_0 +\epsilon } u(a+re^{i \theta })d \theta  < \frac{2 \epsilon u(a)}{2 \pi }.$$  
  于是 
  $$ \begin{aligned}
    u(a) \leq & \frac{1}{2 \pi }\left[  \int_{\theta _0 -\epsilon }^{\theta_0 +\epsilon } + \int_{[0,2 \pi ] \setminus [\theta _0-\epsilon ,\theta _0+\epsilon ]} \right] \cdot u(a+r e^{i \theta }) d \theta \\
    < & \frac{1}{2 \pi }\big( 2 \epsilon u(a)+ (2 \pi -2 \epsilon )u(a)\big)\\
    =& u(a),
  \end{aligned}  $$
矛盾. 

设 $b \in \Omega $, 取 $\gamma $ 连接 $a,b$. 取 $r<<1, a=a_0 \to a_1 \to \cdots \to a_m=b$, s.t. $a_{j+1} \in \Delta (a_j,r)$ 且 $u \in \sh(\Delta (a_j,r))$. 
于是 $u(b)=u(a)=\max_{\Omega }u$.       
\end{proof}

\noindent (7) 肥皂泡原理: 
设 $u \in \usc(\Omega )$, 则 $u \in \sh(\Omega ) \iff \Delta (a,r) \subset \subset \Omega , \forall h \in H(\Delta (a,r)) \cap C(\overline{\Delta (a,r)})$, 若 $u \leq h$ 于 $\partial \Delta (a,r)$, 则 $u \leq h$ 于 $\Delta (a,r)$.      
\begin{proof}
  $\Rightarrow $: $u-h \in \sh (\Delta (a,r))$ 且 $u-h \leq 0$ 于 $\partial \Delta (a,r)$, 所以由 (6) 知 $u-h \leq 0$ 于 $ \Delta (a,r)$;
  
  $\Leftarrow $: 因为 $u \in \usc(\Omega )$, 所以存在 $\{u_j\} \subset C(\overline{\Delta (a,r)}), u_j \searrow u$. 
  记 $P u_j$ 为以 $u_j|_{\partial \Delta (a,r)}$ 为边值得Dirichlet解. 则 $Pu_j|_{\partial \Delta (a,r)}=u_j \geq u$,
  $$ \begin{aligned}
    \Rightarrow \quad& u \leq Pu_j \quad \text{于} ~ \Delta (a,r),\\
     \Rightarrow \quad& u(a) \leq Pu_j(a)=\frac{1}{2 \pi }\int_{0}^{2 \pi } u_j(a+r e^{i \theta }) d \theta .
  \end{aligned}  $$      
  令 $j \to \infty $, 利用Levi定理. 
\end{proof}

\noindent (8) $\sh=\sh_{loc}$.
\begin{proof}
  设 $\Delta (a,r) \subset \subset \Omega , h \in H(\Delta (a,r))\cap C(\overline{\Delta (a,r)})$, s.t. $u \leq h $ 于 $\partial \Delta (a,r)$, $u-h \in \sh_{loc} (\Delta (a,r))$.    
  由 (6) 知 $u-h \leq 0$ 于 $\Delta (a,r)$. 由 (7) $\Rightarrow u \in \sh$.   
\end{proof}
 
\begin{remark}
  $u \in \usc(\Omega )$, 则 $u \in \sh(\Omega )$ $\iff \forall f \in \mathcal{O}(\overline{\Delta (a,r)})$, 若 $u \leq \Re f$ 于 $\partial \Delta (a,r)$, 则    $u \leq \Re f$ 于 $\Delta (a,r)$. \textcolor{red}{(练习) }
\end{remark}

\begin{example}
  $f \in \mathcal{O}(\Omega ) \Rightarrow \log|f|,|f|^{\alpha } \in \sh(\Omega ), \alpha >0$.
  
  \noindent 解: 设 $f \not \equiv 0$, 则 $f^{-1}(0)$ 离散. 则
  $$ u:=\log|f| \in H(\Omega \setminus f^{-1}(0)) .$$ 
  取 $u_j:=\max\{u,-j\}$, 则 $u_j \in \sh(\Omega \setminus f^{-1}(0))$, 且 $\{u<-j\}=\{|f|<e^{-j}\}$ 为 $f^{-1}(0)$ 的邻域. 在其上, $u_j=-j$ 为次调和 $\Rightarrow u_j \in \sh(\Omega )$且 $u_j \searrow u \Rightarrow u \in \sh (\Omega )$. 而 $|f|^{\alpha }=e^{\alpha \log|f|}$, 由 (5) $\Rightarrow |f|^{\alpha } \in \sh(\Omega )$.         
\end{example}

\noindent (9) $u \in H(\Omega ) \iff \pm u \in \sh(\Omega )$.
\begin{proof}
  $\Rightarrow $: 显然. 
  
  $\Leftarrow $: $\pm u \in \usc(\Omega ) \Rightarrow u \in C(\Omega )$. 
  任取 $\Delta (a,r) \subset \subset  \Omega $, 令 $h=Pu$  为 $u|_{\partial \Delta (a,r)}$ 的Dirichlet问题解. 则 $\pm u=\pm h$ 于 $\partial \Delta (a,r)$, 由 (7) 知 $\pm u \leq \pm h$ 于 $\Delta (a,r)$, $\Rightarrow u=h$   于 $\Delta (a,r) \Rightarrow u \in H(\Omega )$.
\end{proof}

\noindent (10) $u \in C^2(\Omega )$. 则 $u \in \sh(\Omega ) \iff \Delta u \geq 0$.  
\begin{proof}
  $\Rightarrow $: 设 $a \in \Omega $, 考虑Taylor 展开 
  $$ \begin{aligned}
    u(a+r e^{i \theta })=&u(a) + \frac{\partial u}{\partial z}(a) \cdot  r e^{i \theta }+ \frac{\partial u}{ \partial \bar z}(a) \cdot  r e^{-i \theta }\\
   &+ \frac{1}{2} \frac{\partial^2 u}{\partial z^2}(a) \cdot ( r e^{i \theta })^2+ \frac{1}{2} \frac{\partial^2 u}{\partial \bar z^2}(a) \cdot ( r e^{-i \theta })^2+ \frac{\partial^2 u}{\partial z \partial \bar z}(a) \cdot  r ^2+o(r^2).
  \end{aligned}  $$
 两边作用 $\frac{1}{2 \pi }\int_{0}^{2 \pi }$ 得: 
$$ \begin{aligned}
  &\frac{1}{2 \pi }\int_{0}^{2 \pi } u(a+r e^{i \theta }) d \theta = u(a)+\frac{\partial^2 u}{\partial z \partial \bar z}(a) \cdot  r ^2+o(r^2)\\
  \Rightarrow \quad & \frac{\partial^2 u}{\partial z \partial \bar z}(a) +o(1) \geq 0 \quad \Rightarrow \quad  \frac{\partial^2 u}{\partial z \partial \bar z}(a) \geq 0 .
\end{aligned}  $$

$\Leftarrow $:  
首先设 $\Delta u >0$. 
设 $\Delta (a,r) \subset \subset \Omega ,h \in H(\Delta (a,r)) \cap  C(\overline{\Delta (a,r)})$, s.t. $u \leq h$ 于 $\partial \Delta (a,r)$. 
令 $v:=u-h$, 则 $v \leq 0$ 于 $\Delta (a,r)$: 若不然, $\exists b \in \Delta (a,r)$, s.t. $v(b)>0$, 则 
$v$ 在 $\Delta (a,r)$ 中某点达到最大值, 于是 $\Delta  \leq 0$ 于该点, $\Rightarrow \Delta u \leq 0$, 矛盾. 于是由 (7) 知 $u \in \sh(\Omega )$. 
一般地, 令 
$$u_j=u+\frac{|z|^2}{j},$$
则 $\Delta u_j >0 \Rightarrow u_j \in \sh (\Omega ) \Rightarrow u \in \sh (\Omega )$, 因为 $u_j \searrow u$. (cf. (4))   
\end{proof}
 
\noindent (11) \textbf{Hardy 定理}: 
设 $u \in \sh (\Delta (0,R)), 0<R \leq + \infty $. 
则 
$$ \phi (r):=\frac{1}{2 \pi } \int_{0}^{2 \pi } u(r e^{i \theta }) d \theta  $$
关于 $\log r \in [-\infty ,\log R)$ 为一个凸增函数 (即 $\phi (e^t)$ 关于 $t$ 凸增).  
\begin{proof}
  令 
  $$ \Phi (z):= \frac{1}{2 \pi } \int_{0}^{2 \pi } u(z e^{i \theta }) d \theta , \quad z \in \cc.$$ 
  下证 $\Phi \in \sh(\Delta (0,R))$: 
  由 Fatou 引理知 $\Phi \in \usc$.  
  设 $\Delta (a,r) \subset \subset \Delta (0,R)$, 则
  $$ \begin{aligned}
     &\frac{1}{2 \pi } \int_{0}^{2 \pi } \Phi  (a+ re^{i \theta }) d \theta\\
    =& \frac{1}{(2 \pi )^2} \int_{0}^{2 \pi } \int_{0}^{2 \pi } u(a e^{i t }+r e^{ i (\theta +t)}) d t d \theta \\
    (\text{Tonelli: 局部有上界})=& \frac{1}{(2 \pi )^2} \int_{0}^{2 \pi } \int_{0}^{2 \pi } u(a e^{i t }+r e^{ i (\theta +t)})  d \theta  d t\\
    (\text{周期性})=& \frac{1}{(2 \pi )^2} \int_{0}^{2 \pi } \int_{0}^{2 \pi } u(a e^{i t }+r e^{ i \theta })  d \theta  d t\\
    (u \in \sh) \geq & \frac{1}{2 \pi } \int_{0}^{2 \pi }  u(a e^{i t })    d t =\Phi (a).
  \end{aligned}  $$ 

  下证 $\phi (r) \to u(0)\, (r \to 0)$: 注意到 
  $$ \Phi (z)= \frac{1}{2 \pi }\int_{0}^{2 \pi } u(|z| \cdot e^{i(\theta +\arg z)}) d \theta =\phi (|z|) ,$$
  由最大模原理, 
  $$ \Phi (z) \leq \max_{|\xi |=r} \Phi (\xi )=\phi (r), \quad \forall |z|<r .$$
  若 $r'=|z|<r \Rightarrow \phi (r') \leq \phi (r)$, 即 $\phi (r) \nearrow $.  
   注意到 
  $$ \phi (r) \geq \phi (0)=u(0), $$ 
  又 $\limsup _{r \to 0} \phi (r) \leq \limsup _{z \to 0} \Phi (z)=\Phi (0)=u(0)$. (因为 $\Phi \in \usc$.) 

  最后证明 $\phi (r)$ 关于 $\log r $ 凸: 设 $0<r_1<r_2<R$. 取 $s,t \in \rr$, s.t. 
  $$ t \log r_j +s =\phi (r_j) .$$    
  定义 
  $$ \psi (z):=\Phi (z)-(t \log |z|+s) \in \sh(\Delta _{r_2}\setminus \overline{\Delta _{r_1}}). $$
  那么 $\psi =0$ 于 $\partial \Delta _{r_1} \cup \partial \Delta _{r_2}$, 由最大模原理, $\psi (z) \leq 0$.   令 $r=|z|$, $ \Rightarrow  \phi (r) \leq t \log r+s$, 即 $\phi (r)$ 关于 $\log r$ 凸.  
  \textcolor{purple}{(设 $\log r=\lambda \log r_1+(1-\lambda )\log r_2$, 则上式等价于 $\phi (r) \leq \lambda \phi (r_1)+(1-\lambda )\phi (r_2)$.) }
\end{proof}

\begin{exercise}
  设 $u \in \sh(\Delta (0,R)), 0<R \leq +\infty $. 
  那么 $\phi (r):=\max_{|z|=r} u(z)$ 关于 $\log r \in [-\infty , \log R)$ 为一个凸增函数. (\textcolor{purple}{(弱)} Hadamard 三圆定理的推广: $M_j:=\max_{\partial \Delta _{r_j}}, r_1<r_2<r_3 \Rightarrow M_2^{\log \frac{r_3}{r_1}}\leq M_1^{\log \frac{r_3}{r_2}} M_3^{\log \frac{r_2}{r_1}}$.)   

  \noindent \textcolor{purple}{(解: cf. Protter和Weinberger-1999-Maximum principles in differential equations. 推出弱三圆定理: 对 $-\log(a-u) \in \sh$ 运用强三圆定理?) }
\end{exercise}

\noindent (12) \textbf{Liouville 定理}: 
 $u \in \sh(\cc), u \leq 0 \Rightarrow u \equiv const$. (\textcolor{red}{可以推出全纯函数的Liouville 定理: 令 $u:=|f|-\max|f| \Rightarrow |f| \equiv const$, 利用开映射定理 $\Rightarrow f \equiv const$. }) 
\begin{proof}
  因为 $\phi (r)$ 关于 $\log r$ 凸增, 且 $\phi \leq 0$, 所以 $\phi (r)=\phi (0)=u(0)$. 则
  $$ \begin{aligned}
    u(a) \leq & \frac{1}{\pi r^2} \int_{\Delta (a,r)} u\\
    (u \leq 0) \leq & \frac{1}{\pi r^2} \int_{\Delta (0, r-|a|)} u\\
    =& \frac{1}{\pi r^2} \int_{0}^{r-|a|} \underbrace{\int_{0}^{2 \pi } u(s e^{i \theta }) d \theta}_{2 \pi  \phi (s)=2 \pi u(0)} s ds\\
    =& u(0) \frac{(r-|a|)^2}{r^2} \to u(0) \quad (r \to \infty ).
  \end{aligned}  $$     
  由 $a$ 的任意性, $u(0) = \max_{\cc}u \Rightarrow u \equiv u(0)$.  
\end{proof}

\begin{exercise}
  设$u \in \sh(\Omega )$ 且 $u(z)=o(\log|z|)~ (z \to \infty )$, 则 $u \equiv const$.   
\end{exercise}

\noindent (13) \textbf{拼接原理}: $\Omega $: 区域, $D \subset \Omega $ 开集, $u \in \sh(\Omega ), v \in \sh(D)$ 且 
$$ \varlimsup _{D \ni z \to \xi } v(z) \leq u(\xi ), \quad \forall \xi \in \partial D \cap \Omega  .$$   
则 
$$ \tilde{u}(z):= \begin{cases}
  u(z) &\text{} z \in \Omega \setminus D\\
  \max \{u(z),v(z)\} &\text{} z \in D
\end{cases}  \quad \in \sh (\Omega ). $$
\begin{remark}
  $\sh$ 函数很柔软. 
\end{remark}
\begin{proof}
  首先, $\tilde{u} \in \usc (\Omega \setminus \partial D)$. 
  设 $a \in \Omega  \cap \partial D, \forall \epsilon >0, \exists \delta >0$, s.t. 
  $$ \begin{aligned}
    u(z) \leq& u(a)+\epsilon , \quad z \in \Delta (a,\delta )\\
   v(z) \leq & \varlimsup_{D \ni w \to a} v(w)+\epsilon  \leq  u(a)+\epsilon, \quad z \in D \cap \Delta (a,\delta ).
  \end{aligned}  $$  
  于是 $\tilde{u} \leq u(a)+\epsilon =\tilde{u}(a)+\epsilon , \forall z \in \Delta (a,\delta )$ $\Rightarrow \tilde{u} \in \usc (\Omega )$.
  
  设 $\Delta (a,r) \subset \subset \Omega $. 不妨设 $\partial D \cap \Delta (a,r) \neq \emptyset $. 
  下面验证肥皂泡原理\textcolor{purple}{(验证次均值性似乎更简单)}. 
  设 $h \in C(\overline{\Delta (a,r)}) \cap H(\Delta (a,r))$, s.t. $\tilde{u} \leq h$ 于 $\partial \Delta (a,r)$.     
  令 $D_a:=D \cap \Delta (a,r)$. 在 $\partial \Delta (a,r) \cap D$ 上, $v \leq \tilde{u} \leq h$; $\partial D \cap  \Delta (a,r) $ 上, $v \leq u \leq h$. 
  因为 $v -h \in \sh(D_a)$, 由最大模原理, 
  $v \leq h$ 于 $D_a$. 于是 $\tilde{u} \leq h$ 于 $D_a$ $\Rightarrow $  $\tilde{u} \leq h$ 于 $\Delta (a,r)$   $\Rightarrow \tilde{u} \in \sh(\Omega )$.      
\end{proof}


\noindent (14) 可去奇性: 
设 $E=v^{-1}(-\infty ) \subset \Omega $ 为闭集, 其中 $v \in \sh (\Omega )$ 且 $v \not \equiv -\infty $.   设 $u \in \sh(\Omega \setminus  E)$ 且在 $\Omega $ 上局部有上界, 则 
$$ \tilde{u}(z):= \begin{cases}
  u(z) &\text{} z \in \Omega \setminus  E\\
  \varlimsup \limits_{\Omega \setminus E \ni \xi \to z} u(\xi )&\text{} z \in E 
\end{cases}  \quad \in \sh(\Omega ) .$$  

\begin{proof}
  由于 $\sh$ 函数为局部性质, 所以不妨设 $u,v <0$.
  $\forall \epsilon >0$, 令 
  $$ u_{\epsilon }(z)= \begin{cases}
    u(z)+\epsilon v(z) &\text{} z \in \Omega \setminus E\\
    - \infty  &\text{} z \in E.
  \end{cases}  $$  

  下证 $u_{\epsilon } \in \sh(\Omega )$: 令 $u_{\epsilon ,j}:=\max\{u_{\epsilon },-j\}$ $\Rightarrow u_{\epsilon ,j} \in \sh(\Omega \setminus E).$
  由于 $\{\epsilon v <-j\}$ 为 $E$ 的邻域, 且其上 $u_{\epsilon ,j}=-j$ 为次调和 $\Rightarrow u_{\epsilon ,j} \in \sh (\Omega ) \Rightarrow u_{\epsilon ,j} \searrow u_{\epsilon } \in \sh(\Omega ).$       

  下证 $\tilde{u}(z) = \big( \sup_{\epsilon >0}u_{\epsilon }(z) \big)^{*}$(cf. (3)):   
  若 $z \in \Omega \setminus E, $ 则
  $$ \sup_{\epsilon >0}u_{\epsilon }(z)=\lim_{\epsilon \to 0}u_{\epsilon }(z)=u(z)=\tilde{u}(z) .$$ 
  由于 $\tilde{u}=u$ 于 $\Omega \setminus E$, 故 $\usc$ $\Rightarrow \tilde{u}(z) = \big( \sup_{\epsilon >0}u_{\epsilon }(z) \big)^{*}$.
  若 $z \in E$, 则 
  $$ \begin{aligned}
    \big( \sup_{\epsilon >0}u_{\epsilon }(z) \big)^{*}=& \varlimsup_{\xi \to z} \sup_{\epsilon >0}u_{\epsilon }(\xi )\\
    =& \varlimsup_{E \not\ni  \xi \to z} \sup_{\epsilon >0}u_{\epsilon }(\xi )\\
    =& \varlimsup_{E \not\ni  \xi \to z} u(\xi )\\
    =& \tilde{u}(z).
  \end{aligned}  $$     
  ~~
\end{proof}

\begin{corollary}
  设 $E=v^{-1}(-\infty )$ 如上. 则 $\Omega \setminus E$ 连通.  
\end{corollary}
\begin{proof}
  假设 $\Omega \setminus E=U_1 \cup U_2$, $U_1,U_2$ 为互不相交非空开集.
  定义 
  $$ u(z):= \begin{cases}
    -1 &\text{} z \in U_1\\
     0 &\text{} z \in U_2
  \end{cases} \quad \in \sh(\Omega \setminus E) .$$  
  因为 $u \leq 0$, 由 (14), 存在 $\tilde{u} \in \sh(\Omega )$ s.t. $\tilde{u}|_{\Omega \setminus E}=u$, 且 $\tilde{u} \leq 0$. 但是 $\tilde{u}|_{U_2}=0$, 由最大模原理 $\Rightarrow \tilde{u} \equiv 0$ 于 $\Omega $, 矛盾.       
\end{proof}


\section{多次调和函数}

\begin{definition}[Oka, Lelong]
  $\Omega \subset \cc^n$ 区域, $u \in \usc(\Omega )$. 若 $\forall a \in \Omega , \forall $ 经过 $a$ 的复线 $l_a$, 有 $u \in \sh(\Omega \cap l_a)$, 则 称 $u$ 为多次调和函数(\textbf{plurisubharmonic functions}), 记为 $\psh (\Omega )$.        
\end{definition}

\begin{remark}
 (i) $v \equiv - \infty $ 可看成是 $\psh$; 
 (ii) 此定义是最适合多复变的定义.
\end{remark}

\begin{remark}
  $u \in \psh (\Omega ) \iff u \in \usc(\Omega )$ 且 $\forall a \in \Omega , \forall \xi \in \cc^{n}, \forall r >0$, 若 $a+\xi  \cdot  \overline{\Delta (0,r)} \subset \Omega $, 则 
  $$ u(a) \leq \frac{1}{2 \pi } \int_{0}^{2 \pi } u(a+\xi \cdot  r e^{i \theta }) d \theta . $$   
\end{remark}

\begin{example}
  设 $u \in \Omega $ 上是凸的(蕴含 $u$  是 Lipschitz 连续的), 则 $u \in \psh (\Omega )$.
  
 \noindent \textbf{解}: 
  $$ \begin{aligned}
    u(a) \leq & \frac{1}{2} [ u(a+\xi \cdot r e^{i \theta })+u(a-\xi \cdot e^{i \theta }) ]\\
   \Rightarrow \quad u(a) \leq & \frac{1}{4 \pi } \int_{0}^{2 \pi } [ u(a+\xi \cdot r e^{i \theta })+u(a-\xi \cdot e^{i \theta }) ] d \theta \\
   (\text{周期性})=&  \frac{1}{2 \pi } \int_{0}^{2 \pi }  u(a+\xi \cdot r e^{i \theta })d \theta .
  \end{aligned}  $$
\end{example}

\begin{example}
 设 $f \in \mathcal{O}(\Omega )$, 则 $\log|f|,|f|^{\alpha }, \alpha >0 \in  \psh(\Omega )$.  
\end{example}

\begin{remark}
  次调和函数的性质, 除(7), (9)-(11) 外, 可平行推广至高维或者从次调和情形推出. 
  例如 (12) Liouville 定理: 设 $u \in \psh(\cc^n ), u \leq 0 \Rightarrow u \equiv const$. 证明: 设 $l $ 为过 $0$ 点的复线, 则 $u|_{l} \in \sh$ 且 $u \leq 0$, 所以 $u \equiv u(0)$ 于 $l$. 且对任意 $z \in \cc^n \setminus \{0\}$ 与 $0$ 点 可由一条复线连接, 所以 $u \equiv u(0)$ 于 $\cc^n$.          
\end{remark}

\begin{remark}
  设 $P=\Pi_{\nu =1}^{n} \Delta (a_{\nu },r_{\nu }), u \in \psh (P)$.  
  则 
  $$ u(a) \leq \frac{1}{|P|} \int_{P}u. $$
  证明: 
  $$ \begin{aligned}
    u(a) \leq & \frac{1}{2 \pi } \int_{0}^{2 \pi }u(a_1+ \xi _1 e^{i \theta },a_2,\cdots , a_{n}) d \theta _1\\
    \leq & \frac{1}{(2 \pi)^2 } \int_{0}^{2 \pi }\int_{0}^{2 \pi }u(a_1+ \xi _1 e^{i \theta },a_2+\xi _2 e^{i \theta }, a_3 \cdots , a_{n}) d \theta _1 d \theta _2\\
    \leq & \frac{1}{(2 \pi)^n } \int_{0}^{2 \pi } \cdots \int_{0}^{2 \pi } u(a_1+ \xi _1 e^{i \theta },\cdots , a_{n}+\xi _n e^{ i \theta }) d \theta _1 \cdots d \theta _n.
  \end{aligned}  $$
  两边作用 $\int_{0}^{r_1} \cdots \int_{0}^{r_n}s_1 \cdots s_n ds_1 \cdots ds_{n}$ 得  
  $$ u(a) \cdot  \frac{r_{1}^{2} \cdots r_{n}^{2}}{2^n}  \leq \frac{1}{(2 \pi )^n} \int_{P}u.$$ 
\end{remark}

\begin{proposition}
  设 $u \in \psh (\Omega ), u\not \equiv - \infty $. 则 $u \in L_{loc}^{1}(\Omega )$.  
\end{proposition}
\begin{proof}
  不妨设 $u <0$. 
  设 $E:=\{z \in \Omega : u ~\text{于 }~ z ~\text{某邻域可积 }\}$. \\
  $E$ 为开集: 显然. \\
  $E$ 非空:   因为 $u \not \equiv - \infty $, 所以 $\exists a \in \Omega $, s.t. $u(a)>-\infty $. 若 $\Delta (a,r) \subset \subset \Omega $, 则 
  $$ \frac{1}{|\Delta (a,r)|^n}\int_{\Delta (a,r)^n}(-u) \leq -u(a) < \infty , \quad \Rightarrow a \in E .$$    
  $E$ 为闭集: 设 $a \in \Omega \setminus E$. 则 $\forall r <<1$, 有 
  $$ \int_{\Delta (a,r)} u =- \infty . $$  
 取  $r << 1,$ s.t. $\forall z \in \Delta (a,r)^n:\Delta (a,r)^n \subset \Delta (z,2r)^n \subset \Delta (a,3r)^n \subset \subset \Omega $. 则 
 $$ \begin{aligned}
  u(z) \leq & \frac{1}{| \Delta (z,2r) |^n } \int_{\Delta (z,2r)^n} u\\
  (u <0)~\leq & \frac{1}{| \Delta (z,2r) |^n } \int_{\Delta (a,r)^n} u \\
  =& -\infty .
 \end{aligned}  $$    
 所以 $\Omega \setminus E$ 为开集. 
 
 因为 $\Omega $ 连通, 所以 $E=\Omega $.  
\end{proof}


\begin{definition}
  设 $E \subset \Omega $, 若 $\forall a \in E$, $\exists U \ni a$, 以及 $u \in \psh(U)$, $u \not \equiv - \infty $, s.t. 
  $$ E \cap U \subset  u^{-1}(-\infty ).$$
  则 称 $E$ 是一个\textbf{多极集(pluripolar set)}. 
  若进一步, 有 $E \cap U=u^{-1}(-\infty )$, 则称 $E$ 为完备多极集.       
\end{definition}

\begin{example}
  解析超曲面为完备多极集.
\end{example}

\begin{corollary}
  多极集为零测集. (\textcolor{red}{反之不对, $1/3$-Cantor 集为零测集, 但非多极集.})
\end{corollary}

\begin{example}
  设 $\{w_n\} \subset \Delta $ 为一个稠密点列, $\{a_n\}$ 为一列正数, s.t. $\sum a_n <\infty $. 令 
  $$ u(z):=\sum_{n}^{\infty } a_n \log |z-w_n|, \quad z \in \cc. $$  
  则 
  \begin{enumerate}[(1)]
    \item $u \in \sh (\cc)$, 且 $u \not \equiv  - \infty $;
    \item   $u^{-1}(-\infty )$ 包含 $\Delta $ 中一个第二纲稠密集;
    \item $u$ 在 $\Delta $ 上 a.e. 不连续.    
  \end{enumerate}
   \begin{proof}
    (1) 设 $R >0$. 当 $|z|<R$ 时,   
    $$ \sum_{j=1}^{n} a_j \log \frac{\left| z-w_j \right| }{R+1} \searrow \sum_{j=1}^{\infty } a_j \log \frac{\left| z-w_j \right|}{R+1} , $$    
    $$ \sum_{j=1}^{n}a_j \log(R+1) \to \sum_{j=1}^{\infty }a_j \log(R+1). $$
   $\Rightarrow u \in \sh (\Delta (0,R)). $  
   由 $R$  的任意性 $\Rightarrow u \in \sh(\cc)$. \\
   若 $\left| z \right| >1 $, 则 $u(z) \geq \sum_{n}^{} \log (\left| z \right| -1 ) >- \infty $.  

   (2) 令 $E=u^{-1}(-\infty )$, 则 
   $$ \{w_n\} \subset E \subset \bar{\Delta }, \quad \bar{\Delta } \setminus E = \cup _{n} \{z \in \bar{\Delta }:u(z) \geq -n\}:= \cup F_n .$$
   因为 $u \in \usc$, 所以 $F_n \subset \bar{\Delta }$ 闭.   
   但 $\{w_j\} \cap F_n =\emptyset , \forall n$, 且 $\{w_j\} \subset \bar{\Delta }$ 稠密, 于是 $F_n$ 疏集. 所以 $\bar{\Delta } \setminus E$ 为第一纲集. 而 由 Baire 定理知, $\bar{\Delta }$ 为第二纲集, 所以 $E=\bar{\Delta }-(\bar{\Delta } \setminus E)$ 为第二纲集. 
   
   (3) $u$ 在 $\bar{E} \setminus E =\bar{\Delta } \setminus E$ 上不连续, 而 $\left| E \right| =0$, 所以 $u$ 在 $\bar{\Delta }$ 上a.e. 不连续.     
   \end{proof}
\end{example}

\begin{remark}
  设 $K \subset \cc$ 为紧集. 则 $K$ 为极集 $\iff$ $K$ 的代数容量为零. (cf. Ransford, Potential Theory.)    
\end{remark}

\begin{proposition}
  设 $E \subset \Omega $ 为闭多极集, 则 $\Omega \setminus E$ 为连通.   
\end{proposition}
\begin{proof}
  (仿照之前推论的证明.)
\end{proof}

\begin{lemma}[Hartogs]\label{lem2.1}
  设 $\{u_j\} \subset \psh(\Omega )$ 内闭一致有上界, 即 
  $$\forall z \in \Omega : \varlimsup_{j \to \infty }u_j(z) \leq M. $$  
  则 $\forall \epsilon >0, \forall $ 紧子集 $K \subset \Omega , \exists j_0=j_0(\epsilon ,K,M) \in \zz^{+}$, s.t. 
  $$ u_j(z) \leq M+\epsilon , \quad \forall z \in K, \forall j \geq j_0. $$  
\end{lemma}
\begin{remark}
  类似于泛函的共鸣定理.
\end{remark}
\begin{proof}
  不妨设 $\Omega $ 有界. 令 $\delta :=d(K, \partial \Omega ), \hat{\Omega }:=\{z \in \Omega : d(z, \partial \Omega ) > \delta /2\}$. 
  则 $K \subset \hat{\Omega } \subset \subset \Omega $. 
  因为 $\{u_j\}$ 在 $\overline{\hat{\Omega }}$ 上一致有界, 所以不妨设 $u_j<0$ 于 $\overline{\hat{\Omega }}$.
  
  令 $\delta '=\delta / \sqrt{n}, r= \delta '/4, 0<r'<r$, s.t. 
  $$ P(z,r) \subset P(w,r+r') \subset \subset \hat{\Omega }, \quad \forall z \in K, w \in P(z,r') .$$ 
  由于 
  $$ \varlimsup_{j \to \infty } \fint_{P(z,r)} u_j \leq  \fint_{P(z,r)} \varlimsup_{j \to \infty } u_j \leq M, $$
  所以 
  $$ \exists j_0 >0, ~s.t.~ \fint_{P(z,r)} u_j \leq M +\frac{\epsilon }{2}, \quad \forall j \geq j_0 .$$
  设 $w \in P(z,r')$, 则 
  $$ \begin{aligned}
    u_j(w) \leq& \fint_{P(w,r+r')} u_j \leq \frac{\left| P(z,r) \right| }{\left|P(w,r+r')\right|} \fint_{P(z,r)} u_j \\
    \leq& \left(  \frac{r}{r+r'} \right)^{2n} \left( M+\frac{\epsilon }{2} \right) \\
    \leq& M+ \epsilon. \quad (\text{若} r'=r'(\epsilon ,M) <<1.)
  \end{aligned}  $$
使用 Heine-Borel 定理即可.
\end{proof}

\begin{remark}
  $u \in \psh(\Omega ), u \not \equiv -\infty  \Rightarrow u \in BMO_{loc}(\Omega )$, 即 
  $$ \forall \Omega ' \subset \subset \Omega \Rightarrow  \sup_{B \subset \Omega '} \fint_B \left| u-u_B \right| <\infty , \quad u_B:=\fint_B u. $$ 
\end{remark}


\begin{definition}
  设 $u \in C^2 (\Omega )$. 称 
  $$ \mathcal{L}u(a;\xi ):= \sum_{j,k=1}^{n}\frac{\partial^2 u}{\partial z_j \partial \bar{z}_k}(a) \xi _j \bar{\xi }_k $$ 
  为 $u$ 的Levi形式; $\left( \frac{\partial^2 u}{\partial z_j \partial \bar{z}_k} \right)$ 为 $u$ 的复Hessian矩阵.  
\end{definition}

\begin{proposition}
  设 $u \in C^2(\Omega ; \rr)$, 则 
  $$ u \in \psh (\Omega ) \iff \mathcal{L} u (a;\xi ) \geq 0, \quad \forall a \in \Omega , \forall \xi \in \cc^n ,$$ 
  即 $\left( \frac{\partial^2 u}{\partial z_j \partial \bar{z}_k} \right)$ 半正定 ( $i \partial \bar{\partial }u \geq 0$). 
\end{proposition}
\begin{remark}
  $i dz \wedge d \bar{z}>0$. 
\end{remark}
\begin{proof}
  考虑复线 $z=a+t \xi , t \in \cc$, 则 
  $$ \frac{\partial^2 u }{\partial t \partial \bar{t}}(a+t \xi )=\mathcal{L}u(a;\xi ). $$
  ~~ 
\end{proof}

\begin{definition}
  设 $u \in C^2(\Omega ;\rr)$. 若 $\mathcal{L}u(a;\xi )>0, \forall a \in \Omega , \forall \xi \in \cc^n \setminus \{0\}$ (即 $\left( \frac{\partial^2 u }{\partial z_j \partial \bar{z}_k} \right)>0$ 或 $i \partial \bar{\partial }u>0$), 则称 $u$ 在 $\Omega $ 上强多次调和, 其全体记为 $\spsh (\Omega )$.     
\end{definition}

\begin{remark}
  $u \in \spsh (\Omega ) \Rightarrow i \partial \bar{\partial } u$ 为 $\Omega $ 上一个 K\"ahler 度量.  
  反之, 若 $\omega $ 为$\Omega $ 上一个 K\"ahler 度量, 则局部地, $\omega = i \partial \bar{\partial } u, u \in \spsh (\Omega )$.  
\end{remark}

\begin{theorem}\label{th2.1}
  设 $u \in \psh(\Omega ), \forall \epsilon >0$, 令 $\Omega _{\epsilon }:=\{z \in \Omega : d(z, \partial  \Omega )>\epsilon \}$. 则存在 $u_{\epsilon } \in \spsh (\Omega_{\epsilon } )$, s.t. $u_{\epsilon } \searrow u$. 
  若进一步假设 $u \in C(\Omega )$, 则 $u_{\epsilon }$ 内闭匀敛于 $u$.       
\end{theorem}

\begin{definition}[磨光算子]
  $\chi \in C_{0}^{\infty }(\cc^n), \chi \geq 0, \supp \chi \subset \overline{\mathbb{B}^n}$, 且 
  $$\int_{\cc^n} \chi =1, \quad \chi (z_1,\cdots ,z_n)=\chi (|z_1|,\cdots ,|z_n|).$$  
\end{definition}
\begin{example}
  $$ \chi (z) = \begin{cases}
    \frac{1}{\gamma } e^{\frac{1}{\left| z \right|^2-1 }}, &\text{若}~ |z| <1\\
    0, &\text{若}~ |z| \geq 1
  \end{cases} , \quad \gamma =\int_{\mathbb{B}^n}  e^{\frac{1}{\left| z \right|^2-1 }}.$$
\end{example}

设 $u \in L_{loc}^{1}(\Omega ), \forall  \epsilon >0$, 令 
$$ \chi _{\epsilon }(z)=\frac{1}{\epsilon ^{2n}} \chi \left( \frac{z}{\epsilon } \right), \quad z \in \Omega _{\epsilon };$$ 
$$ u_{\epsilon }(z):=u * \chi _{\epsilon }(z)=\int_{\xi \in \cc^n} u(\xi ) \chi _{\epsilon } (z-\xi )=\int_{\xi \in \cc^n } u(z-\xi \cdot \epsilon )\chi (\xi ).$$

\begin{lemma}
  (1) 若$u \in C(\Omega )$, 则 $u_{\epsilon } \to u$ 内闭一致;\\
  (2) 若 $u \in \psh (\Omega ) $, 则$ u_{\epsilon } \in \psh (\Omega _{\epsilon })$.  
\end{lemma}
\begin{proof}
  (1) 设 $K \subset \Omega $ 紧. 
  令
  $$ \omega (\epsilon ):=\sup_{z \in K, \xi \in \overline{\mathbb{B}^n}}\left| u(z-\epsilon \xi )-u(z) \right| \to 0, ~ (\epsilon \to 0). $$  
  则
  $$ \left| u_{\epsilon }(z)-u(z) \right| \leq \int_{\xi \in \cc^n} \left| u(z-\epsilon \xi )-u(z) \right| \chi (\xi ) \leq \omega (\epsilon ) \to 0, \quad \forall  z \in K. $$

  (2) 设 $a+\xi \cdot  \overline{\Delta (a,r)} \subset \Omega _{\epsilon }$, 则 
  $$ \begin{aligned}
    \frac{1}{2 \pi }\int_{0}^{2 \pi } u_{\epsilon }(a +\xi r e^{i \theta}) d \theta =& \frac{1}{2 \pi }\int_{0}^{2 \pi } \int_{\omega  \in \cc^n} u(a+\xi r e^{i \theta }-\epsilon \omega )\chi (\omega )dv_{\omega }d \theta \\
  (\text{Fubini})  =& \int_{\cc^n} \chi (\omega ) dv_{\omega } \cdot  \frac{1}{2 \pi } \int_{0}^{2 \pi } u(a+\xi r e^{i \theta }-\epsilon \omega )d \theta \\
  \geq & \int_{\cc^n} u(a-\epsilon \omega  ) \chi (\omega ) dv_{\omega }\\
  =& u_{\epsilon }(a).
  \end{aligned}  $$ 
  ~~
\end{proof}

\begin{proof}[定理 \ref{th2.1}]
  先取 $u_{\epsilon }$ 为 $u$ 的磨光化. 下证  
  $u_{\epsilon } \searrow u$: 令 $\delta (z)=d(z,\partial \Omega )$, 
  则 $z \in \Omega _{|t|}, \forall |t| <\delta (z)$. 
  令 
  $$ g(t):= \int_{\cc^n} u(z-t \xi ) \chi (\xi ) = \frac{1}{|t|^{2n}}\int_{\cc^n} u(\xi ) \chi (\frac{z-\xi }{t}).$$
  于是 $g(t)=g(|t|)=u_{|t|}(z)$. 
  利用与 Hardy 定理类似的证明 $\Rightarrow g$ 次调和.
  由最大模原理, $u_{\epsilon }(z) =g(\epsilon ) \searrow $.
  设 $u(z) \neq  - \infty $. 因为 $g$ 上半连续, 
  $$ u(z)=g(0) \leq g(\epsilon ) < g(0)+\eta , \quad(\eta <<1), $$
  因此 $u_{\epsilon }(z) \searrow u(z)$. 
  若 $u(z) = -\infty $ (\textcolor{red}{练习}).
  
  最后, 将 $u_{\epsilon }+\epsilon |z|^2$ 代替 $u_{\epsilon }$ 即可.  
\end{proof}

\begin{corollary}
  设$F \in \mathcal{O}(\Omega ,\Omega '), u \in \psh(\Omega ')$. 则 $F^{*}u=u \circ F \in \psh (\Omega )$.  特别地, 多次调和性为双全纯不变量($\Rightarrow $ 可推广到复流形上).
\end{corollary}
\begin{proof}
  先设 $u \in C^2(\Omega ')$. 令 $w=F(z)$, 则
  $$ \frac{\partial^2 u \circ F}{\partial z_j \partial \bar{z}_k} = \sum_{\alpha \beta }^{} \frac{\partial^2 u}{\partial w_{\alpha } \partial \bar{w}_{\beta }} \cdot  \frac{\partial w_{\alpha }}{\partial z_j} \cdot  \frac{\partial \bar{w}_{\beta }}{\partial \bar{z}_k}.$$
  于是 
  $$ \begin{aligned}
    \mathcal{L}u \circ F (a;\xi )=& \sum_{j,k} \frac{\partial^2 u \circ F}{\partial z_j \partial \bar{z}_k} (a) \xi _j \bar{\xi }_k \\
    =& \sum_{\alpha \beta } \sum_{j,k} \frac{\partial^2 u}{\partial w_{\alpha } \partial \bar{w}_{\beta }} \cdot  \frac{\partial w_{\alpha }}{\partial z_j} \cdot  \overline{\frac{\partial w_{\beta }}{\partial z_k}}  \xi _j \bar{\xi }_k ~
    \geq 0.
  \end{aligned}  $$  

  一般地, 取 $u_{\epsilon } \in \spsh(\Omega_{\epsilon }^{\prime } ), u_{\epsilon }  \searrow u$.  则 $u_{\epsilon } \circ F \in \psh(F^{-1}(\Omega _{\epsilon }^{\prime }))$, 且 $\cup _{\epsilon } F^{-1}(\Omega _{\epsilon }^{\prime })=\Omega, u_{\epsilon } \circ F \searrow u \circ F \Rightarrow u \circ F \in \psh (\Omega )$.  
\end{proof}

\begin{corollary}
  设 $\lambda$  为 $\rr^m$  上一个 $C^2$  凸函数, 且关于每个分量 $\nearrow $. 设$u_1,\cdots ,u_m \in \psh(\Omega )$,  则 
  $ \lambda (u_1,\cdots ,u_m) \in \psh(\Omega ) .$
\end{corollary}
\begin{proof}
  先设 $u_j \in C^2(\Omega )$.
  $$ \Rightarrow ~ \frac{\partial^2 \lambda (u_1,\cdots u_m) }{\partial z_j \partial \bar{z}_k} = \sum_{\alpha }^{}\frac{\partial \lambda }{\partial x_{\alpha }} \frac{\partial^2 u_{\alpha } }{\partial z_j \partial \bar{z}_k}+\sum_{\alpha ,\beta } \frac{\partial^2 \lambda  }{\partial x_{\alpha } \partial x_{\beta }} \frac{\partial u_{\beta }}{\partial z_j} \frac{\partial u_{\alpha }}{\partial \bar{z}_k}.$$ 

  一般地, 取 $u_j,\epsilon \in \spsh(\Omega _{\epsilon })$, s.t. $u_{j,\epsilon } \searrow u_j$. 则 $\lambda (u_{1,\epsilon},\cdots u_{m,\epsilon } ) \in \spsh (\Omega _{\epsilon })$, 且 $\lambda (u_{1,\epsilon},\cdots, u_{m,\epsilon } ) \searrow \lambda (u_1,\cdots ,u_m) \Rightarrow \lambda (u_{1},\cdots, u_{m } )\in \psh(\Omega )$.    
\end{proof}

\begin{example}
  设 $f_1, \cdots ,f_m \in \mathcal{O}(\Omega ), \alpha _1, \cdots ,\alpha _m >0$. 则 
  $$ u(z):=\log(|f_1|^{\alpha _1}+\cdots +|f_m|^{\alpha _m}) \in \psh(\Omega ).$$
  特别地, \underline{解析子集为完备多极集}. 
\end{example}
\begin{proof}
  令 $\lambda (x_1,\cdots ,x_m):=\log(e^{x_1}+ \cdots +e^{x_m})$. 则 
  $$ u(z)=\lambda (\alpha _1 \log |f_1|,\cdots ,\alpha _m \log|f_m|) .$$ 
  只要证明 $\lambda $ 凸且关于每个变量递增即可. 
  后者显然. 
  因为
  $$ \frac{\partial^2 \lambda  }{\partial x_{\alpha } \partial x_{\beta }}=\frac{\partial }{\partial x_{\alpha }}\frac{e^{x_{\beta }}}{e^{x_1}+ \cdots +e^{x_m}}= \frac{e^{x_{\beta }}\delta _{\alpha \beta }}{e^{x_1}+ \cdots +e^{x_m}}-\frac{e^{x_{\beta }}e^{x_{\alpha }}}{(e^{x_1}+ \cdots +e^{x_m})^2},$$ 
  所以 $\forall \xi \in \rr^m$, 
  $$ \begin{aligned}
   \sum_{\alpha ,\beta }^{} \frac{\partial^2 \lambda  }{\partial x_{\alpha } \partial x_{\beta }} \xi _{\alpha }\xi _{\beta }=& \frac{\sum_{\alpha }^{} e^{x_{\alpha}}(\xi _{\alpha })^2}{e^{x_1}+ \cdots +e^{x_m}}-\frac{\sum_{\alpha }^{}(e^{x_{\alpha }} \xi _{\alpha })^2}{(e^{x_1}+ \cdots +e^{x_m})^2}\\
    =& \frac{\sum_{\alpha }^{} e^{x_{\alpha}}(\xi _{\alpha })^2 \sum_{\alpha }^{} e^{x_{\alpha }}-\sum_{\alpha }^{}(e^{x_{\alpha }} \xi _{\alpha })^2}{\left(\sum_{\alpha }^{} e^{x_{\alpha }} \right)^2}\\
    \geq & 0 \quad (\text{Schwarz 不等式}).
  \end{aligned}  $$ 
\end{proof}

\noindent \textbf{问题}: 定理 \ref{th2.1} 中是否可取 $u_{\epsilon }\in \spsh(\Omega )$? (一般不行, 但 $\Omega $ 为拟凸域可以.) 


\section{多次调和函数的应用}

\begin{theorem}[Hartogs]
  设 $f$ 为 $\Omega $ 上的一个函数, 其关于每一个分量全纯. 则 $f \in \mathcal{O}(\Omega )$.   
\end{theorem}
\begin{proof}
  Step 1. 假设命题 $n-1$ 时成立 ($n=1$ 是平凡的). 设 $\Delta _1,\dots,\Delta _n$  为圆盘, s.t. $\bar{\Delta }_1 \times \dots \times \bar{\Delta }_n \subset \Omega $, 则存在圆盘 $\widetilde{\Delta }_n \subset \Delta _n$, s.t. $f \in L^{\infty }(\Delta _1 \times \cdots \times  \widetilde{\Delta }_n)$.    
  证明如下: 对$\forall m \in \zz^{+}$, 令 
  $$E_m := \{z_n \in \bar{\Delta }_n: \forall z' \in \bar{\Delta }_1 \times \cdots \times \bar{\Delta }_{n-1}, \left| f(z',z_n) \right| \leq m\}.$$  
  则 $E_m= \cap _{z' \in \bar{\Delta }_1 \times \cdots \times \bar{\Delta }_{n-1}} \{z_n \in \bar\Delta _n: \left| f(z',z_n) \right| \leq m\}$ 为闭集. 且 $\bar{\Delta }_n= \cup _m E_m$ (归纳假设). 由 Baire 纲定理, 存在 $m$ s.t. $E_{m}^{\circ } \neq \emptyset $, 只要取 $\widetilde{\Delta }_n \subset E_{m}^{\circ }$即可.  
  
  Step 2. 设 $P(a,R) \subset \subset \Omega $, 且 $\exists 0<r<R$, s.t. $f \in \mathcal{O}(P(a',R) \times \Delta (a_n,r))$. 则 $f \in \mathcal{O}(P(a,R))$. 证明如下: 不妨设 $a=0$, 且 $\left| f \right| \leq M$ 于 $P(0',R) \times \Delta (0,r)$.
  $$ f(z',z_n)=\sum_{j}c_j(z')z_{n}^{j}, \quad z \in P(0,R), \quad (*)$$
  其中 $c_j(z')=\frac{1}{j!} \frac{\partial^j f}{\partial z_{n}^{j}}(z',0) \in \mathcal{O}(P(0',R))$.
  由 Cauchy 估计 $\Rightarrow |c_j(z')| \leq M/r^j$.    令 
  $$ u_j(z'):=\frac{1}{j} \log \left| c_j(z') \right| \in \psh(P(0',R)). $$     
  任取 $0<R_1<R_2<R$, 因为 $(*)$ 收敛, 所以 
  $$ c_j(z')R_{2}^{j} \to 0, \quad \forall z' \in P(0',R_2). $$
  所以 
  $$ \varlimsup_j u_j(z') \leq -\log R_2 .$$
  由 Hartogs 引理 (引理 \ref{lem2.1}), 
  $$ u_j(z) \leq -\log R_1, \quad \forall z' \in P(0',R_1), \forall j \geq j_0 >> 1. $$ 
  $\Rightarrow (*)$ 在 $P(0',R_1)$ 内闭匀敛 \textcolor{purple}{(确保连续性)} $\Rightarrow f \in \mathcal{O}(P(0,R_1)) \Rightarrow f \in \mathcal{O}(P(0,R))$.  

  Step 3. 设 $a \in \Omega $, 取 $P(a,2R) \subset \Omega $. 由 Step 1, 存在 $\Delta (c_n,r) \subset \Delta (a_n,R)$, s.t. $f \in L^{\infty }(P(a',R) \times \Delta (c_n,r))$; 由 Step 2 可知 $f \in \mathcal{O}(P((a',c_n),R))$ \textcolor{purple}{(回顾Osgood 定理)}. 因为 $a \in P((a',c_n),R)$, 所以 $f \in \mathcal{O}_a$.       
\end{proof}


\noindent\textbf{Dirichlet 问题}: 设 $\Omega \subset \subset \cc^n, f \in C(\partial \Omega ,\rr)$. Perron 方法: 
$$ \mathcal{P}:=\{u \in \psh (\Omega ):\varlimsup_{z \to \xi }u(z) \leq f(\xi ), \forall \xi \in \partial \Omega \} $$
令 $u_f:=\sup_{\mathcal{P}}u$, ---Perron 函数. 
$u_f$ 是哪一类函数? (正则性, etc.).  $u_f|_{\partial \Omega }=f$? 

\begin{remark}
  可导出 $\left( i \partial \bar{\partial }  \right)^n u_f=0$ (分布意义). 
  $$ \partial \bar{\partial } f \wedge \cdots \wedge \partial \bar{\partial } f =0 \quad \text{复齐次Monge-Ampere 方程}.$$
  若 $u_f \in C^2$, 则 
  $$ \left( i \partial \bar{\partial }  \right)^n u_f= \det \left( \frac{\partial^2 u_f }{\partial u_j \partial \bar{u}_k} \right) \cdot dV, \quad n \geq 2. $$ 
  若 $\Omega =\mathbb{B}^n$, 则 $u_f \in C^{1,1}(\mathbb{B}^n)$, 且 $u_f|_{\partial \mathbb{B}^n}=f$.   
\end{remark}
\begin{itemize}
   \item $\bar{\partial }$-方程;
   \item $\partial \bar{\partial }$-方程;
   \item M-A 方程. (Ref: Blocki: Complex M-A operator.)  
\end{itemize}


\begin{definition}
  设 $u$ 为 $\Omega $ 上实值函数. 若 $\pm u \in \psh(\Omega )$, 则称 $u$ 在 $\Omega $ 上\textbf{多调和(pluriharmonic)}. 记全体为 $\ph (\Omega )$.     
\end{definition}

\begin{proposition}
  (1) $u \in \ph (\Omega ) \iff u \in C^{\infty }(\Omega )$ 且 $\dbb u=0$.\\
  (2) $u \in \ph (\Omega ) \iff u$ 可以局部表示为一个全纯函数的实部.  
\end{proposition}
\begin{proof}
  (1) $\Leftarrow $: 显然. \\
  $\Rightarrow $: 首先, $u$ 在过 $\Omega $ 中任一点 $a$ 的复线上为调和的.   若 $P(a,r) \subset \Omega $, 则 
  $$ \begin{aligned}
    u(z)=& \frac{1}{2 \pi }\int_{0}^{2 \pi } \frac{r^2-\left| z_1-a_1 \right| ^2}{\left| r e^{i \theta _1}-(z_1-a_1) \right| ^2} u(a_1+r e^{i \theta _1},z_2,\dots,z_n) d \theta _1\\
    =& \cdots \\
    =&  \frac{1}{(2 \pi )^n}\int_{0}^{2 \pi } \cdots \int_{0}^{2 \pi } \Pi _{j=1}^{n} \frac{r^2-\left| z_j-a_j \right| ^2}{\left| r e^{i \theta _j}-(z_j-a_j) \right| ^2} u(a_1+r e^{i \theta _1},\dots,a_n+r e^{i \theta _n}) d \theta _1 \cdots d \theta _n
  \end{aligned}  $$   
  $\Rightarrow u \in C^{\infty }(\Omega )$.
  在复线 $z=a+t\xi, t \in \cc $ 上, 
  $$ 0=\frac{\partial^2 u (z+t \xi ) }{\partial t \partial \bar{t}} =\mathcal{L}u (z;\xi ).$$  
  $\Rightarrow \dbb u =0$.
  
  (2) $\Leftarrow $: 若局部地, $u =\Re (f), f \in \mathcal{O}$. 则 
  $$ \dbb u =\frac{1}{2} \dbb (f+\bar{f})=0, $$
  由 (1) 即得.  \\
  $\Rightarrow $: 任取球 $B(a,r ) \subset \subset \Omega $. 
  因为 $\partial u$ 为 $B(a,r)$ 的一个可微 $1$-形式, 所以
  $$ d(\partial u)= \bar{\partial } \partial u=-\dbb u =0. $$     
  由 Poincar\'e 引理知, 存在 $f \in C^{\infty }(B(a,r))$, s.t. $df=\partial u$. 
  由于 
  $$ d f =\partial f +\bar{\partial }f ,$$
  所以 $\db f =0$, i.e. $f \in \mathcal{O}(B(a,r))$. 
  由 
  $$ \begin{aligned}
    &d(u-2\Re f)= \partial u+\db u -df-d \bar{f}=0\\
    \Rightarrow~ & u-2 \Re f \equiv const.\\
    \Rightarrow~ &u=\Re (2f+const).
  \end{aligned}  $$  
  ~~
\end{proof}


\begin{theorem}[Riemann可去奇性定理 III]
  设 $E \subset \Omega $ 为闭多极集, $f \in \mathcal{O}(\Omega \setminus E) \cap L_{loc}^{\infty }(\Omega )$. 则 $\exists !~ \tilde{f} \in \mathcal{O}(\Omega )$, s.t. $\tilde{f}|_{\Omega \setminus E}=f$.    
\end{theorem}
\begin{proof}
  因为 $\pm \Re f \in \ph (\Omega \setminus E)$, 且于 $\Omega $ 局部有界, 所以由多次调和可去奇性定理 $\Rightarrow \pm \Re f$ 可延拓为 $u_{\pm}\in \psh (\Omega )$. 显然, 
  $$ u_{+}+u_{-}=0, \quad \text{于} ~ \Omega \setminus E. $$
  设 $a \in E$, 取 $P(a,r) \subset \Omega $. 
  $$ (u_{+}+u_{-})(a) \leq \fint_{P(a,r)} (u_{+}+u_{-})=0 \quad (\text{多极集Lebsgue测度为0}). $$      
  另一方面, 
  $$ (u_{+}+u_{-})(a) \leq \varlimsup_{\Omega \setminus E \ni z \to a} (u_{+}+u_{-})(z)=0. $$
  所以 $u_{+}=-u_{-}$ 于 $\Omega $ $\Rightarrow  u_{\pm} \in \ph(\Omega )$.   

  同理, $\pm \Im f$ 也可延拓为 $\pm v \in \ph(\Omega )$. 令 
  $$ \tilde{f}:=u_{+}+i v_{+} ,$$
  则 $\tilde{f}|_{\Omega \setminus E}=f$.
  于是 $\db \tilde{f}=0$ 于 $\Omega \setminus E$. 由于 $\Omega \setminus E$ 在 $\Omega $ 中稠密, 且 $\tilde{f} \in C^{\infty }(\Omega )$, 所以 
  $$ \db \tilde{f}=0 \quad \text{于}~ \Omega , $$    
  即 $\tilde{f} \in \mathcal{O}(\Omega )$.     
\end{proof}


\begin{corollary}[Rado]
  若 $f \in C(\Omega ) \cup \mathcal{O}(\Omega \setminus E)$,  其中 $E=f^{-1}(0)$.  则 $f \in \mathcal{O}(\Omega )$.  
\end{corollary}

\begin{proof}
  $u:=\log |f| \in \usc (\Omega ) \cap \psh (\Omega \setminus E)$. 因为 
  $$ u_j:=\max \{u,-j\} \in \psh (\Omega \setminus E) ,$$
  且在 $E$  的某邻域内 $=-j \in \psh $, 所以  
  $u_j \in \psh(\Omega )$, 且 $u_j \searrow u$, 所以 $u \in \psh (\Omega )$.
  于是 $E=f^{-1}(0)=u^{-1}(-\infty) $ 为一个闭多极集. 由Riemann可去奇性III, $\exists \tilde{f} \in \mathcal{O}(\Omega )$, s.t. $\tilde{f}|_{\Omega \setminus E}=f$. 但是 $\Omega \setminus E$ 于 $\Omega $ 中周密且 $f \in C(\Omega )$, 所以 $\tilde{f}=f$ 于 $\Omega $.           
\end{proof}

\begin{corollary}[Osgood]
  设 $\Omega _1, \Omega _2$ 为 $\cc^n$ 中有界区域, $F \in \mathcal{O}(\Omega _1,\Omega _2)$ 且为双射. 则 $F$ 为双全纯.    
\end{corollary}
\begin{proof}
  (\textcolor{red}{提示: 应用秩定理 + Riemann 可去奇性定理 III}).
\end{proof}
